\hypertarget{abstract}{%
\section{Abstract}\label{abstract}}

Existing cascade systems like BranchyNet, MSDNet, EdgeBERT, RouteLLM use
hardcoded confidence thresholds or hardware timers to route requests.
They do not use live metrics from infrastructure state. This study
explores a context-aware agentic gateway which uses an edge context
protocol, a dynamic Medallion pipeline, MCP skills, PPO agents,
bidirectional loop and a policy orchestration layer to make dynamic
routing decisions. Here the PPO is trained offline in a Gymnasium
simulation with domain randomization which generates a frozen policy.
This policy is deployed to the PPO agent to make routing decisions. For
steady and bursty types of traffic, the proposed system reduces cloud
escalations from 100\%, i.e. static-threshold baselines, down to 19.5\%
under both, eliminating SLA violations under bursty load entirely (down
from 65--66\% for static thresholds) at a small, consistent accuracy
cost (about one percentage point). This result required two rounds of
diagnosis and correction during evaluation: the first trained policy was
undertrained and was silently governed by its own deterministic safety
fallback rather than its learned action, and once fixed, an
unclamped-latency reward term was found to give the policy no gradient
against escalating once an SLA was already violated, which it
disproportionately did under sustained network degradation. Both are
reported as part of the evaluation, not concealed. The system was
compared against five baselines and eight literature baselines across
three traffic scenarios, plus a seven-run ablation, including a
governance-layer ablation comparing a shared frozen policy against a
policy orchestration layer that can reuse, fine-tune, or train a new
policy per client and, since a small local LLM was added as an
oversight mechanism, escalate its own decision when warranted.

